%% =================================================================
%% Wei, Mei, Zhao, Xiao, Sun, Zhang, Guo.
%% "Nondestructively identifying the mechanical behavior of soft
%%  tissues using surface deformation with an explicit inverse
%%  approach."
%% Appl. Math. Modelling 134 (2024) 126--147.
%%
%% Reconstruction of page 8 (printed p. 133) as study notes.
%% Prose: extracted verbatim from the PDF text layer via pdftotext.
%% Math:  hand-transcribed from the rendered page image —
%%        Eqs. (27)–(28), the final two gradient expressions from
%%        the inclusion-embedded case.
%% Figure 6: placeholder box with caption + visual description.
%% NOTE: this page begins the Results section.
%% =================================================================

\documentclass[11pt]{article}

\usepackage[margin=1in]{geometry}
\usepackage{amsmath, amssymb}
\usepackage{mathrsfs}
\usepackage{bm}
\usepackage{soul}
\definecolor{hlyellow}{RGB}{255,255,170}
\definecolor{hlcyan}{RGB}{170,255,255}
\definecolor{hlgreen}{RGB}{170,255,170}
\definecolor{hlpink}{RGB}{255,170,170}
\definecolor{hlorange}{RGB}{255,221,170}
\usepackage{fancyhdr}

\pagestyle{fancy}
\fancyhf{}
\lhead{\small Y.~Mei et al.}
\rhead{\small Applied Mathematical Modelling 134 (2024) 126--147}
\cfoot{\small \thepage}
\renewcommand{\headrulewidth}{0.4pt}

\setcounter{page}{133}
\setcounter{equation}{26}  % next equation is (27)

\begin{document}

\begin{equation}
\frac{\partial\mu}{\partial\mu^{0}}
= \prod_{i=1}^{n} H\!\left(\phi^{i}(\mathbf{x})\right)
\end{equation}

For the shear modulus $\mu^{j}$ of the $j$-th component,
$\frac{\partial\mu}{\partial\mu^{j}}$ is given by:
\begin{equation}
\frac{\partial\mu}{\partial\mu^{j}}
= 1 - H\!\left(\phi^{j}(\mathbf{x})\right)
\end{equation}

\section*{Results}

\subsection*{\textit{Case 1: A layered rectangular structure}}

\noindent
In the first numerical example, we analyze a layered rectangular
structure, illustrated in Fig.~6. The dimensions of this structure
are 1.0\,mm $\times$ 0.2\,mm. To discretize the problem domain, we
utilize 7200 linear triangular elements. In the forward finite
element simulation, a concentrated force is applied to the top edge,
while the bottom edge was restricted the vertical motion and the
center point is fully fixed. In the inverse problem, our focus lies
in utilizing displacement datasets collected from the top edge as
inputs for the inverse scheme. Acknowledging that relying on a single
surface displacement dataset may not produce optimal reconstruction
results, we change the location of the concentrated force and
incorporate multiple surface displacement datasets to solve the
inverse problem. The target shear modulus distributions for one (the
homogenous case), two, and three layers are depicted in Fig.~6(a),
(b), and (c), respectively. The shear modulus values of each layer in
the structure are of a similar order to that of skin tissue [36].
Considering the uncertainty regarding the total number of layers, we
initially assume a configuration with six layers for the inverse
problem, as depicted in Fig.~7(a). For the explicit inverse scheme,
we utilize 10 control points to interpolate each interface. Thus, the
total number of optimization variables is 56 including 50 geometric
variables and 6 shear moduli.

From Fig.~7, it can be seen that both the shape of the interfaces
between the neighboring layers and the shear modulus values of each
layer varies and finally reach to the target shear modulus
distribution (Sample 1). \pdfcomment[icon=Comment,color=cyan]{User input: how fast is the algorithm in the Wei paper?\\\\Claude comments: The paper reports speed ONLY in iteration counts, NEVER in wall-clock time (seconds/minutes). No "CPU time", "runtime", or "elapsed time" appears anywhere in the 22-page paper. This is a significant omission — iteration counts are not directly comparable to actual speed because: (1) each iteration's cost depends on problem size and FEM solve cost, (2) explicit and implicit methods may have different per-iteration costs, (3) hardware/implementation details are never disclosed.\\\\The specific numbers available across the paper: - Layered case (page 8): 52 iterations for explicit, \~{}10x less than implicit - Inclusion case (page 17-18): 328 iterations for explicit vs 4354 for   implicit (13.3x speedup in iterations)\\\\Without wall-clock timing, "how fast" cannot be answered in absolute terms (e.g., "runs in 5 minutes"). The paper only answers "how many fewer iterations" — which is a relative measure, not an absolute one.}\sethlcolor{hlcyan}\hl{Note that it takes only 52 iterations for
convergence.} Figs.~8--10 are the reconstructed shear modulus
distributions for Samples 1--3 using explicit and implicit inverse
approaches for the noise free cases, respectively. In the implicit
inverse method, the total number of optimization variables is
linearly proportional to the mesh size. In this case, the total
number of optimization variables is 3691, which is equal to the total
number of nodes. The reconstructed results demonstrate that both
inverse methods can yield decent reconstruction with noise free
cases. Further, the implicit method fails to recover the
heterogeneity at the upper corner compared with the results by the
explicit inverse method. \pdfcomment[icon=Comment,color=cyan]{User input: how fast is the algorithm in the Wei paper?\\\\Claude comments: The paper reports speed ONLY in iteration counts, NEVER in wall-clock time (seconds/minutes). No "CPU time", "runtime", or "elapsed time" appears anywhere in the 22-page paper. This is a significant omission — iteration counts are not directly comparable to actual speed because: (1) each iteration's cost depends on problem size and FEM solve cost, (2) explicit and implicit methods may have different per-iteration costs, (3) hardware/implementation details are never disclosed.\\\\The specific numbers available across the paper: - Layered case (page 8): 52 iterations for explicit, \~{}10x less than implicit - Inclusion case (page 17-18): 328 iterations for explicit vs 4354 for   implicit (13.3x speedup in iterations)\\\\Without wall-clock timing, "how fast" cannot be answered in absolute terms (e.g., "runs in 5 minutes"). The paper only answers "how many fewer iterations" — which is a relative measure, not an absolute one.}\sethlcolor{hlcyan}\hl{Since the total number of optimization
variables is remarkably reduced in the explicit inverse method, the
total number of minimization iterations is reduced. The total number
of minimization iterations for the explicit inverse method is
approximately 10 times less than that required for the implicit
inverse method.} Furthermore, the shear modulus reconstruction
obtained through the implicit

% ---------- Figure 6 placeholder ----------
\vspace{1em}
\begin{center}
\fbox{\parbox{0.92\textwidth}{\centering\vspace{1em}
\textbf{[Figure 6 — placeholder]}\\[0.8em]
Three layered rectangular structures (1.0\,mm $\times$ 0.2\,mm),
each showing the target shear modulus (SM) distribution in MPa
with a jet-style colorbar on the right.\\[0.4em]
\textbf{(a) Sample 1}: homogeneous, single layer, SM $= 3.00$\,MPa
(uniform blue). Downward arrows on top edge indicate concentrated
force locations. Red text ``displacement measurement'' with
leftward arrows on left side.\\[0.3em]
\textbf{(b) Sample 2}: two layers. Bottom layer SM $\approx 3$\,MPa
(blue), top layer SM $\approx 10$\,MPa (red). Green circles on top
boundary, red circles on bottom boundary mark displacement
measurement points.\\[0.3em]
\textbf{(c) Sample 3}: three layers (blue--red--blue sandwich). SM
ranges 1.00--10.00\,MPa. Same boundary marker pattern as (b).\\[0.3em]
All three samples share the same mesh (7200 triangular elements) and
boundary conditions (top-edge force, bottom-edge vertical
constraint, center point fixed).
\vspace{1em}}}
\end{center}
\vspace{0.3em}
\begin{center}
\textbf{Fig. 6.} The target shear modulus distribution of three
samples with different number of layers (Unit: MPa). SM stands for
shear modulus.
\end{center}

\end{document}
